\paragraph{A one-sided inverse-square shielding budget.}
\label{par:one-sided-shielding}
The integrated requirement \eqref{eq:integratedshieldu}
does not need two-sided control of a whole neighborhood
of the consecutive pair. There is a sharper geometric
comparison that penalizes \emph{only outward gaps larger
than the gap under protection}. This identifies a
particular observable of the moving zero configuration
whose heat-time integral would suffice.

Suppose $a<b$ are consecutive real simple zeros,
$g=b-a>0$, and enumerate the exterior real zeros
as $a-\ell_m$ and $b+r_m$ $(m\ge1)$, with
$\ell_m,r_m>0$ increasing. Define successive
exterior gaps
\[
d_{L,1}=\ell_1,\quad
d_{L,j}=\ell_j-\ell_{j-1}\ (j\ge2),
\qquad
d_{R,1}=r_1,\quad
d_{R,j}=r_j-r_{j-1}\ (j\ge2).
\]
In particular
$\ell_m=\sum_{j=1}^m d_{L,j}$
and similarly on the right.

\begin{theorem}[Sharp one-sided shielding bound]
\label{thm:onesidedshield}
Assume the bi-infinite real simple zero sequence
has absolutely convergent shielding sum
\[
\mathcal S=\sum_{m\ge1}
\left\{\frac1{\ell_m(\ell_m+g)}
+\frac1{r_m(r_m+g)}\right\}.
\]
Let $D=2-g^2\mathcal S$. Then, for every integer $R\ge1$,
\begin{equation}\label{eq:onesidedfinite}
\boxed{\begin{aligned}
[D]_+ \le{}&\frac2{R+1}\\
&+\sum_{j=1}^R
\left(\frac1{j^2}-\frac1{(R+1)^2}\right)
\frac{
(d_{L,j}-g)_++(d_{R,j}-g)_+}{g}.
\end{aligned}}
\end{equation}
Consequently, in the extended nonnegative real sense,
\begin{equation}\label{eq:onesidedinfinite}
\boxed{
[D]_+\le\frac1g\sum_{j\ge1}
\frac{(d_{L,j}-g)_++(d_{R,j}-g)_+}{j^2}.
}
\end{equation}
These inequalities require no restriction on
outward gaps \emph{smaller} than $g$.
The coefficient $j^{-2}$ is optimal to
first order in a single outward-gap perturbation.
\end{theorem}

\begin{proof}
Set $f(v)=1/(v(v+1))$ for $v>0$.
The two normalized shielding contributions
at index $m$ are $f(\ell_m/g)$ and $f(r_m/g)$.
The regular-lattice reference satisfies
$\sum_{m\ge1}f(m)=1$.
Since
\[
-f'(v)=\frac1{v^2}-\frac1{(v+1)^2}
\]
is positive and strictly decreasing,
the mean-value theorem gives, for $v>0$,
\begin{equation}\label{eq:tangent}
f(m)-f(v)\le
\left(\frac1{m^2}-\frac1{(m+1)^2}\right)
(v-m)_+.
\end{equation}
For $v\le m$ the left side is nonpositive;
for $v>m$ integrate $-f'$ between $m$ and $v$.

Omitting all shielding terms with $m>R$ can
only \emph{increase} $D$. Therefore
\begin{align*}
D
&\le 2-\sum_{m=1}^R
\{f(\ell_m/g)+f(r_m/g)\}\\
&=\frac2{R+1}+
\sum_{m=1}^R
\{f(m)-f(\ell_m/g)+f(m)-f(r_m/g)\}.
\end{align*}
Let $\eta_{L,j}=d_{L,j}/g-1$ and similarly
$\eta_{R,j}$. Then
\[
(\ell_m/g-m)_+
=\left(\sum_{j=1}^m\eta_{L,j}\right)_+
\le\sum_{j=1}^m(\eta_{L,j})_+.
\]
Apply \eqref{eq:tangent} on both sides,
interchange two finite nonnegative sums, and use
the exact telescoping identity
\[
\sum_{m=j}^R
\left(\frac1{m^2}-\frac1{(m+1)^2}\right)
=\frac1{j^2}-\frac1{(R+1)^2}.
\]
This proves \eqref{eq:onesidedfinite};
its right side is nonnegative and therefore
also bounds $[D]_+$.

If the infinite series on the right of
\eqref{eq:onesidedinfinite} diverges there
is nothing further to prove. Otherwise
the finite sums converge to it as $R\to\infty$,
and the first term tends to zero; this proves
the infinite bound.

To see first-order sharpness,
start with a perfect lattice of gap $g$
and enlarge only the $j$th exterior gap on
one side from $g$ to $(1+\delta)g$,
$\delta>0$. Then
\[
D(\delta)=\sum_{m=j}^\infty
\{f(m)-f(m+\delta)\},\qquad
D'(0+)=\sum_{m=j}^\infty
\left(\frac1{m^2}-\frac1{(m+1)^2}\right)
=\frac1{j^2}.
\]
At $\delta=1$ the difference telescopes again:
\[
D(1)=\frac1{j(j+1)}.
\]
Thus the linear inverse-square coefficient
cannot be uniformly lowered even in
an arbitrarily small single-gap perturbation.
\end{proof}

\begin{corollary}[Conditional one-sided
heat-time collision barrier]\label{cor:onesidedheat}
Along any backward heat interval on which the
zero-motion equation \eqref{eq:shieldgap} holds,
let $d_{L,j}(s),d_{R,j}(s)$ denote the moving
exterior gaps. For a fixed integer $R\ge1$ put
\[
\mathcal E_R(s)=\frac2{R+1}
+\sum_{j=1}^R
\left(\frac1{j^2}-\frac1{(R+1)^2}\right)
\frac{
(d_{L,j}(s)-g(s))_+
+(d_{R,j}(s)-g(s))_+}{g(s)}.
\]
Then
\begin{equation}\label{eq:one-sided-dynamic}
G(s)\ge G(0)-4\int_0^s\mathcal E_R(v)\,dv.
\end{equation}
In particular, if
\begin{equation}\label{eq:one-sided-budget}
\boxed{
4\int_0^\sigma\mathcal E_R(v)\,dv
<G(0)
}
\end{equation}
throughout the real-simple evolution up to
a proposed terminal time $\sigma$,
this chosen pair cannot collide by that time.
\end{corollary}
\begin{proof}
By Theorem~\ref{thm:onesidedshield},
$[D(s)]_+\le\mathcal E_R(s)$.
Since $G'=-4D\ge-4[D]_+$,
integration gives \eqref{eq:one-sided-dynamic}.
The strict bound \eqref{eq:one-sided-budget}
forces the gap square to stay positive.
\end{proof}

The improvement is structural, not yet a theorem
about actual Riemann zero rigidity. The earlier
near-lattice criterion constrained both signs of
every exterior displacement. Here only the
\emph{positive excess} of each successive exterior
gap enters, and the weights are inverse-square.
The finite-$R$ remainder $2/(R+1)$ is explicit;
over a half-unit backward interval it alone
contributes $4/(R+1)$ to the total budget.
Thus high zeros with $g(0)\asymp1/\log X$
still call for $R$ on the scale of $(\log X)^2$,
even if all weighted excesses vanished.

The single-gap example also exposes why
zero-density information alone is insufficient:
doubling one exterior gap at index $j$
changes the local shielding deficit by
$1/[j(j+1)]$, despite leaving the rest of
the ideal infinite lattice unchanged.
Moreover, the positive-part budget
\eqref{eq:one-sided-budget} is
\emph{sufficient but not necessary}:
a squared gap may remain positive despite
repeated decreases and recoveries whose total
downward variation exceeds its initial value.

The true missing estimate is now explicit:
establish \eqref{eq:one-sided-budget}, or a
comparably effective signed-deficit integral,
for \emph{all} genuine theta-flow zero pairs
through backward time $1/2$, with a further
argument precluding a first loss of real
rootedness at unbounded height. Neither
assertion follows from the current theta
moment bounds, the local seven-unit
Laguerre interval, or the known average
zero-density formula. Consequently
Theorem~\ref{thm:onesidedshield}
does not establish RH.

\paragraph{Exact verification.}
The standalone rational checker
\path{certificates/xi_one_sided_shielding_exact.py}
verifies the scalar tangent identity,
1,800 adversarial finite-gap configurations,
the optimal first-order $j^{-2}$ weights
and the single doubled-gap example.
The Lean module
\path{XiKernel/OneSidedShielding.lean}
certifies the scalar rational
identities and conditional integrated
margin. The genuine infinite zero-motion
equation, interchange of limits and the
all-time theta-specific budget are
analytical steps, not end-to-end
kernel-certified Lean results.
